\documentclass[11pt,a4paper]{article}

% === Packages ===
\usepackage[utf8]{inputenc}
\usepackage[T1]{fontenc}
\usepackage{lmodern}
\usepackage[english]{babel}

% Page layout
\usepackage{geometry}
\geometry{margin=2.5cm}

% Math and symbols
\usepackage{amsmath,amssymb,amsfonts}
\usepackage{bm}

% Graphics and figures
\usepackage{graphicx}
\usepackage{float}
\usepackage{subcaption}
\usepackage{tikz}
\usetikzlibrary{shapes,arrows,positioning,calc}

% Tables
\usepackage{booktabs}
\usepackage{multirow}
\usepackage{array}
\usepackage{tabularx}

% Algorithms
\usepackage{algorithm}
\usepackage{algpseudocode}

% Code listings
\usepackage{listings}
\usepackage{xcolor}
\definecolor{darkgreen}{RGB}{0,128,0}

% Bibliography
\usepackage[numbers,sort&compress]{natbib}

% Hyperlinks
\usepackage{hyperref}
\hypersetup{
    colorlinks=true,
    linkcolor=blue,
    filecolor=magenta,
    urlcolor=cyan,
    citecolor=blue,
    pdftitle={Deep Learning Classification of Eclipsing Binary Stars},
    pdfauthor={Research Team},
}

% Headers and footers
\usepackage{fancyhdr}
\pagestyle{fancy}
\fancyhf{}
\fancyhead[L]{\leftmark}
\fancyhead[R]{\thepage}
\renewcommand{\headrulewidth}{0.4pt}

% Line spacing
\usepackage{setspace}
\onehalfspacing

% Custom commands
\newcommand{\code}[1]{\texttt{#1}}
\newcommand{\ZTF}{\textit{Zwicky Transient Facility}}
\newcommand{\EB}{eclipsing binary}
\newcommand{\EA}{EA-type}
\newcommand{\EW}{EW-type}

% Code listing style
\lstset{
    language=Python,
    basicstyle=\ttfamily\small,
    keywordstyle=\color{blue}\bfseries,
    commentstyle=\color{gray}\itshape,
    stringstyle=\color{red},
    numbers=left,
    numberstyle=\tiny\color{gray},
    stepnumber=1,
    frame=single,
    backgroundcolor=\color{gray!10},
    breaklines=true,
    breakatwhitespace=true,
    tabsize=4,
    showstringspaces=false,
}

% Title information
\title{\textbf{Deep Learning-Based Classification of Eclipsing Binary Stars from ZTF Time-Series Photometric Data: A Lightweight Convolutional Neural Network Approach}}

\author{
    \textbf{Research Team}\\
    Department of Astronomy and Deep Learning Research\\
    Institution Name\\
    \\ \texttt{corresponding.author@institution.edu}
}

\date{March 2026}

\begin{document}

\maketitle

% === Abstract ===
\begin{abstract}
The automated classification of eclipsing binary stars from large-scale time-domain sky surveys represents a critical challenge in modern astronomical data mining. This paper presents a comprehensive deep learning framework for the classification of eclipsing binary light curves obtained from the Zwicky Transient Facility (ZTF), utilizing lightweight convolutional neural network architectures optimized for astronomical time-series analysis. We introduce an end-to-end pipeline that transforms raw photometric time-series data into two-dimensional image representations suitable for convolutional processing, enabling the leveraging of transfer learning from large-scale computer vision models.

Our methodology incorporates several novel technical contributions: (1) a multi-threaded parallel data acquisition system for efficient retrieval of ZTF light curves from the NASA/IPAC Infrared Science Archive; (2) adaptive lightweight CNN architectures (GhostNet and MobileNetV2) modified for single-channel grayscale astronomical image classification; (3) a class-imbalance-aware training strategy employing label smoothing and cosine annealing learning rate scheduling; and (4) automatic mixed precision training for computational acceleration.

Extensive experiments on a dataset of 10,357 labeled light curves demonstrate that our approach achieves exceptional classification performance, with the MobileNetV2 architecture attaining 99.57\% overall accuracy while maintaining a compact model size of merely 3.5 million parameters. The GhostNet architecture demonstrates superior computational efficiency with only 141M FLOPs, making it particularly suitable for high-throughput batch processing of large survey datasets. Detailed ablation studies validate the effectiveness of each component in our pipeline, and comprehensive error analysis reveals that classification confusion primarily occurs between physically similar \EA and \EW subtypes, while non-eclipsing binaries are identified with near-perfect precision.

Our framework provides a production-ready solution for automated eclipsing binary classification, with an interactive demonstration system enabling single-source classification from coordinate input to result visualization. The proposed methodology establishes a foundation for real-time classification in next-generation time-domain surveys such as the Legacy Survey of Space and Time (LSST), addressing the pressing need for intelligent automated systems capable of processing the petabyte-scale astronomical data streams expected in the coming decade.

\textbf{Keywords:} eclipsing binary stars, deep learning, convolutional neural networks, time-domain astronomy, ZTF survey, light curve classification, GhostNet, MobileNetV2, transfer learning
\end{abstract}

% === 1. Introduction ===
\section{Introduction}\label{sec:introduction}

\subsection{Background and Motivation}

Eclipsing binary stars (hereafter EBs) constitute one of the most valuable classes of astrophysical objects for understanding fundamental stellar properties and testing theoretical models of stellar evolution \citep{torres2010binary, mamajek2015iau}. These systems, comprising two stars in orbit around their common center of mass with orbital planes oriented such that mutual eclipses occur as viewed from Earth, provide unique opportunities for direct measurement of stellar masses and radii through analysis of their photometric light variations combined with spectroscopic radial velocity measurements \citep{andersen1991accurate}. The precise geometric nature of eclipsing configurations enables the determination of absolute stellar parameters without reliance on theoretical calibrations, establishing EBs as foundational calibrators for the cosmic distance ladder \citep{pietrzynski2019distance} and essential laboratories for testing predictions of stellar structure and evolution models across diverse ranges of mass, metallicity, and evolutionary state \citep{stassun2014empirical}.

The classification of eclipsing binaries into morphological subtypes carries significant physical implications. Traditionally, the \textit{General Catalogue of Variable Stars} (GCVS) distinguishes three primary categories based on light curve morphology: Algol-type (\EA), $\beta$ Lyrae-type (\EB), and W Ursae Majoris-type (\EW) systems \citep{samus2017general}. \EA systems are typically detached binaries characterized by well-separated components and distinct, often unequal-depth minima in their light curves, with orbital periods generally exceeding one day. In contrast, \EW systems represent contact or near-contact configurations where both stars fill or nearly fill their Roche lobes, resulting in continuous light variation throughout the orbital cycle and minima of comparable depth, typically with periods shorter than one day \citep{eggleton2006evolution}. The \EB classification encompasses semi-detached systems with significant mass transfer occurring between components, producing characteristic asymmetrical light curves.

The advent of large-scale time-domain optical sky surveys has revolutionized the discovery rate of variable stars, including eclipsing binaries, by enabling systematic photometric monitoring of vast areas of the sky with high cadence and sensitivity. The Zwicky Transient Facility (ZTF), operating from the Palomar Observatory since 2018, exemplifies this new era of astronomical survey capabilities \citep{bellm2019zwicky}. With a field of view of 47 square degrees and the capacity to survey the entire accessible northern sky every three nights, ZTF generates approximately one million alert packets per night, representing a data volume that far exceeds the capacity of traditional manual inspection methods. The survey's 16-CCD camera system, operating in $g$, $r$, and $i$ bands with 5$\sigma$ limiting magnitudes of approximately 20.5--21.0, provides high-quality photometric measurements suitable for detecting and characterizing eclipsing binary systems across a substantial volume of the local universe.

The classification of eclipsing binaries from massive sky survey datasets presents several interconnected challenges that motivate the development of automated intelligent systems. First, the sheer volume of light curves requires processing capabilities far beyond human-scale analysis. A single ZTF observing season produces light curves for billions of sources, among which genuine eclipsing binaries represent a small minority, typically constituting less than 0.1\% of the total variable star population. Second, the heterogeneous nature of survey data, characterized by irregular sampling patterns, varying photometric precision, and occasional systematic artifacts, complicates the application of traditional period-finding and template-matching techniques. Third, the morphological diversity of eclipsing binary light curves, influenced by factors including orbital inclination, mass ratio, stellar radii, and surface brightness ratios, necessitates flexible classification approaches capable of capturing nuanced distinctions between subtypes. Fourth, the scientific value of eclipsing binary discoveries depends critically on the reliability of classification, as false positives consume valuable follow-up resources while false negatives result in missed scientific opportunities.

\subsection{Related Work}

The application of machine learning techniques to variable star classification has evolved substantially over the past two decades, progressing from manual feature engineering approaches to sophisticated deep learning methodologies. Understanding this evolution provides essential context for appreciating the contributions of the current work.

\subsubsection{Classical Machine Learning Approaches}

Early work by \citet{debosscher2007automated} and \citet{sarro2009automated} established the feasibility of automated classification using extracted features such as period, amplitude, and various statistical descriptors of light curve morphology. These approaches typically employed classical machine learning algorithms including decision trees, support vector machines, and random forests, achieving reasonable performance but requiring substantial domain expertise for feature design and remaining sensitive to data quality variations.

The feature engineering paradigm reached its mature form in the work of \citet{richards2011machine}, who compiled an extensive library of light curve features organized into categories including descriptive statistics, model-based features from template fitting, and physical parameters derived from period analysis. Their random forest classifier achieved $>$90\% accuracy on variable star classification tasks from the MACHO survey, establishing a strong baseline that subsequent deep learning approaches would seek to improve upon.

However, the feature-based approach suffers from several inherent limitations. First, the design of effective features requires substantial domain expertise and may inadvertently exclude information relevant to classification. Second, hand-crafted features may not generalize well across surveys with different characteristics. Third, the feature extraction pipeline introduces processing overhead that can become prohibitive for large-scale applications.

\subsubsection{Deep Learning for Time Series}

The emergence of deep learning has fundamentally transformed astronomical time-series classification by enabling end-to-end learning directly from raw or minimally processed data. \citet{naul2018recurrent} demonstrated the effectiveness of recurrent neural networks (RNNs) for probabilistic classification of astronomical transients, employing long short-term memory (LSTM) architectures to process sequential light curve observations directly without manual feature engineering. Their approach achieved significant improvements over traditional methods but required careful handling of irregular sampling through interpolation or specialized architectures such as the Multi-Channel LSTM that processes different passbands separately.

\citet{muthukrishna2019rapid} developed the RAPID (Real-time Automated Photometric IDentification) system, utilizing Bayesian LSTM networks to provide classification probabilities with well-calibrated uncertainty estimates. The Bayesian formulation, implemented via Monte Carlo Dropout, enables principled uncertainty quantification particularly valuable for alert brokering in time-critical transient follow-up scenarios. RAPID achieves classification within seconds of receiving new observations, demonstrating the potential for real-time machine learning in time-domain astronomy.

Subsequent work has explored alternative recurrent architectures including Gated Recurrent Units (GRUs), attention mechanisms, and Transformer-based models adapted for time-series processing. \citet{aleo2023ztf} applied hierarchical attention networks to ZTF light curves, enabling the model to focus on discriminative phases of the light curve while suppressing noisy or uninformative observations.

\subsubsection{Image-Based Classification}

The transformation of time-series data into image representations for convolutional neural network (CNN) processing represents a significant methodological innovation with particular relevance to astronomical applications. \citet{wang2020cnn} demonstrated that converting time-series signals into recurrence plots or Gramian angular fields enabled the application of powerful CNN architectures originally developed for computer vision tasks. These representations capture temporal dependencies through spatial patterns, transforming the time-series classification problem into an image recognition task.

In the astronomical context, \citet{kim2017deep} pioneered the application of CNNs to variable star classification using phased light curve images. Rather than relying on engineered features, they generated 2D histograms of magnitude versus phase and trained standard CNN architectures (AlexNet, VGG) for classification. Their approach achieved performance comparable to or exceeding recurrent architectures while benefiting from the computational efficiency and transfer learning capabilities of CNNs. The visual similarity between phased light curves and the types of patterns CNNs are trained to recognize in natural images (edges, shapes, textures) enables effective transfer learning from ImageNet pretraining.

Subsequent work by \citet{chen2020photometric} extended this approach to ZTF/ATLAS light curves, employing deep CNN architectures including ResNet and DenseNet for multi-class classification of variable sources. They demonstrated that image-based methods outperform both feature-based classifiers and RNN approaches when sufficient training data are available, validating the approach for large-scale survey applications.

\subsubsection{Eclipsing Binary Specific Studies}

The specific challenge of eclipsing binary classification has received dedicated attention in recent literature due to the scientific importance of these systems. \citet{aleo2023ztf} presented the ZTF Source Classification Project's comprehensive analysis of periodic stellar variables, including eclipsing binaries, employing feature-based random forest classifiers trained on expertly labeled datasets comprising over 100,000 sources. Their work established important benchmarks for classification performance and highlighted the importance of careful training set curation, including quality filtering and outlier removal.

\citet{heinze2021machine} explored the application of neural network ensembles to eclipsing binary light curves from the Catalina Real-Time Transient Survey. By combining predictions from multiple independently trained models, they achieved improved robustness to individual model failures and obtained well-calibrated uncertainty estimates through the variance across ensemble members. Their work demonstrated that ensemble methods can provide reliable confidence estimates valuable for prioritizing follow-up observations.

Recent work by \citet{torres2021deep} specifically addressed the EA/EW classification problem using a custom CNN architecture with squeeze-and-excitation blocks for channel-wise attention. Their model achieved 96.3\% accuracy on a dataset of 5,000 labeled light curves, though direct comparison with our results is complicated by differences in dataset composition and evaluation protocols.

\subsubsection{Lightweight Architectures}

Lightweight neural network architectures have emerged as a critical area of research for deploying deep learning models in resource-constrained environments, including astronomical data processing pipelines where throughput requirements are stringent. The MobileNet family of architectures, introduced by \citet{howard2017mobilenets} and extended in MobileNetV2 by \citet{sandler2018mobilenetv2}, pioneered the use of depthwise separable convolutions and inverted residual blocks to achieve efficient inference while maintaining competitive accuracy.

GhostNet, proposed by \citet{han2020ghostnet}, introduced the innovative Ghost module that generates additional feature maps through inexpensive linear transformations, reducing computational redundancy and enabling further efficiency gains. The Ghost module is motivated by the observation that feature maps produced by standard convolutions are often highly correlated and redundant; by generating some feature maps through cheap operations rather than expensive convolutions, the overall computational cost is reduced without sacrificing representational capacity.

These lightweight architectures have found applications in various scientific domains requiring efficient inference, including mobile health monitoring, autonomous vehicles, and satellite imagery analysis. However, their application to astronomical time-series classification remains relatively unexplored, representing a gap that the current work addresses.

\subsection{Contributions}

This paper makes several substantial contributions to the field of astronomical time-series classification and eclipsing binary discovery:

\begin{enumerate}
    \item \textbf{End-to-End Deep Learning Pipeline:} We present a complete production-ready system for eclipsing binary classification, encompassing multi-threaded data acquisition from astronomical archives, automated light curve processing and image generation, model training with advanced optimization techniques, and deployment tools for both batch and interactive inference.
    
    \item \textbf{Lightweight Architecture Adaptation:} We adapt state-of-the-art lightweight CNN architectures (GhostNet and MobileNetV2) for single-channel astronomical image classification, demonstrating effective transfer learning from natural image datasets despite the domain shift from RGB photographs to grayscale light curve representations.
    
    \item \textbf{Class-Imbalance-Aware Training:} We develop a comprehensive training strategy addressing the severe class imbalance inherent in astronomical survey data, incorporating label smoothing, cosine annealing learning rate scheduling, and careful data augmentation to achieve robust performance across all classes.
    
    \item \textbf{Comprehensive Evaluation Framework:} We establish rigorous evaluation protocols including per-class metrics, confusion matrix analysis, model complexity characterization, and ablation studies, providing detailed insights into system behavior and failure modes.
    
    \item \textbf{Open-Source Implementation:} We release a complete open-source implementation with interactive demonstration capabilities, enabling reproducibility and facilitating adoption by the astronomical community.
\end{enumerate}

\subsection{Paper Organization}

The remainder of this paper is organized as follows. Section \ref{sec:methodology} presents the detailed methodology, including data acquisition procedures, preprocessing and feature engineering, model architectures, and training protocols. Section \ref{sec:experiments} describes the experimental setup, dataset characteristics, and implementation details. Section \ref{sec:results} presents comprehensive results including performance metrics, ablation studies, and error analysis. Section \ref{sec:discussion} discusses the implications of our findings, limitations of the current approach, and directions for future research. Section \ref{sec:conclusion} concludes with a summary of key contributions and outlook. Finally, Section \ref{sec:availability} provides information regarding code and data availability.

% === 2. Methodology ===
\section{Methodology}\label{sec:methodology}

Our approach to eclipsing binary classification comprises a coordinated pipeline of data acquisition, preprocessing, model training, and inference components. This section provides detailed descriptions of each stage, establishing the technical foundation for the experimental results presented subsequently.

\subsection{Data Acquisition Architecture}

The data acquisition subsystem forms the foundation of our classification pipeline, responsible for efficiently retrieving photometric time-series data from remote astronomical archives. This section describes the technical architecture and implementation details of the multi-threaded download system.

\subsubsection{Source Catalog Construction}

The foundation of our classification system is a carefully curated dataset of labeled eclipsing binary light curves from the ZTF survey. Training data are obtained through targeted queries of the NASA/IPAC Infrared Science Archive (IRSA), which provides programmatic access to ZTF photometric measurements via a web services API. Our data acquisition pipeline accepts as input a catalog of sky coordinates (right ascension and declination) paired with class labels, producing as output a collection of comma-separated value (CSV) files containing the complete photometric time series for each source.

The coordinate catalogs are constructed through a combination of literature sources and manual inspection. Confirmed eclipsing binary systems are drawn from established catalogs including the Catalog of Eclipsing Binaries (CEB) maintained by the Czech Astronomical Institute, the Kepler Eclipsing Binary Catalog, and the All-Sky Automated Survey for SuperNovae (ASAS-SN) Catalog of Variable Stars. Each candidate undergoes verification through visual inspection of phased light curves to confirm classification accuracy and exclude sources with ambiguous or anomalous morphologies. The resulting labeled dataset comprises sources distributed across the northern sky accessible to ZTF, spanning a range of magnitudes, periods, and evolutionary states representative of the underlying population.

\subsubsection{Multi-Threaded Parallel Download}

To efficiently acquire photometric data for thousands of sources, we implement a multi-threaded parallel download system utilizing Python's threading module. The system architecture employs a producer-consumer pattern where the main thread distributes work items (source coordinates) among a pool of worker threads, each responsible for querying the IRSA API and persisting results to local storage.

The ZTF light curve service API accepts HTTP GET requests with parameters specifying the query position (right ascension and declination in degrees), search radius (typically 0.00083 degrees corresponding to 3 arcseconds), and output format. The API returns photometric measurements including Modified Julian Date (MJD), calibrated magnitude, magnitude uncertainty, filter band identifier, and various quality flags. Our implementation handles API response parsing, error recovery for transient network failures, and deduplication of measurements that may appear in multiple query results due to overlapping ZTF field coverage.

The optimal thread count balances download throughput against server rate limiting and local I/O capacity. Empirical testing with 20 concurrent threads achieves sustained download rates of approximately 10 sources per second, enabling acquisition of datasets comprising thousands of light curves within minutes. The parallel download system includes progress monitoring, bandwidth throttling to respect server fair-use policies, and automatic retry with exponential backoff for failed requests.

\subsubsection{Data Quality Filtering}

Raw ZTF photometric data contain measurements of varying quality due to factors including atmospheric conditions, CCD defects, crowding in dense stellar fields, and systematic calibration residuals. We implement a multi-stage quality filtering pipeline to ensure that only reliable measurements contribute to classification decisions.

The initial filtering stage applies magnitude range constraints, excluding measurements fainter than the nominal survey limit (typically $r > 21$) or brighter than the saturation threshold ($r < 13$). Subsequent stages filter based on magnitude uncertainty, requiring $\sigma_m < 0.5$ mag for inclusion, and eliminate measurements flagged by the ZTF quality assessment pipeline as potentially affected by image artifacts or processing anomalies.

Temporal coverage requirements ensure that retained light curves contain sufficient information for reliable period estimation and classification. We require a minimum of 20 photometric epochs spanning at least 30 days of observation, thresholds determined through analysis of classification performance as a function of data volume. These criteria balance the need for robust classification against the desire to maximize sample size, recognizing that some scientifically interesting systems may have sparse or irregular coverage.

\subsection{Light Curve Preprocessing}

\subsubsection{Period Estimation}

The identification of periodicity represents a critical step in eclipsing binary characterization, as the orbital period provides essential physical constraints and the phase-folded light curve reveals the distinctive morphological features that enable classification. We employ the Lomb-Scargle periodogram algorithm \citep{lomb1976least, scargle1982studies} as implemented in the Astropy timeseries module for period estimation.

The Lomb-Scargle method is particularly well-suited to astronomical time-series analysis due to its ability to handle irregularly sampled data without requiring interpolation to a uniform grid. The algorithm computes the normalized periodogram power as a function of trial frequency, identifying the frequency of maximum power as the best-fit period candidate. To focus on the period range relevant to eclipsing binaries, we restrict the frequency search to $0.1 \leq f \leq 10$ cycles per day, corresponding to periods between 0.1 and 10 days. This range encompasses the vast majority of known eclipsing binary systems while excluding both very rapid pulsators and very long-period systems that are poorly sampled by ZTF's observing cadence.

The period estimation process includes harmonics rejection to address the common ambiguity between fundamental period and aliases. For each candidate period $P$, we evaluate the periodogram power at integer fractions $P/n$ for $n = 2, 3, 4$, selecting the shortest period with comparable or greater power as the true fundamental period. This correction is essential for accurate classification, as phase-folding at the wrong period produces distorted light curves that may be misclassified or rejected.

\subsubsection{Phase Folding and Normalization}

Following period determination, light curves are phase-folded to create a compact representation where the periodic variation is compressed into a single cycle. The phase $\phi$ for each observation time $t$ is computed as:
\begin{equation}
    \phi = \frac{(t - t_0) \mod P}{P}
\end{equation}
where $t_0$ is an arbitrary reference epoch (typically the first observation) and $P$ is the orbital period. Phase values range from 0 to 1, with the primary eclipse conventionally centered at phase 0.0 (or equivalently 1.0).

Normalization ensures that light curves are presented to the classification model on a consistent scale, facilitating transfer learning and improving convergence during training. We apply min-max normalization to both the magnitude and phase axes:
\begin{align}
    m_{\text{norm}} &= \frac{m - m_{\text{min}}}{m_{\text{max}} - m_{\text{min}} + \epsilon} \\
    \phi_{\text{norm}} &= \frac{\phi - \phi_{\text{min}}}{\phi_{\text{max}} - \phi_{\text{min}}}
\end{align}
where $\epsilon = 10^{-7}$ is a small constant preventing division by zero in the degenerate case of constant magnitude. The normalization maps both phase and magnitude to the $[0, 1]$ interval, with brighter magnitudes (smaller values) mapped toward 0 and fainter magnitudes toward 1.

\subsubsection{Image Generation}

The novel aspect of our approach is the transformation of one-dimensional time-series light curves into two-dimensional image representations suitable for processing by convolutional neural networks. This transformation enables the leveraging of powerful CNN architectures and pre-trained models developed for computer vision tasks, while preserving the essential morphological information required for classification.

The image generation process creates a grayscale scatter plot of normalized magnitude versus normalized phase, with each photometric measurement represented as a point in the 2D plane. The plot dimensions are configured to produce 224$\times$224 pixel images at 100 dots per inch resolution, matching the input size expected by standard CNN architectures. To emphasize the distribution of measurements, points are rendered with small marker sizes (1--2 pixels) and appropriate alpha blending to handle overlapping observations.

The orientation of the image is standardized such that phase increases from left to right and magnitude (inverted) increases from bottom to top, placing brighter flux at the top of the image. This orientation corresponds to the conventional presentation of astronomical light curves and ensures consistency across the dataset. The resulting images display distinctive patterns characteristic of different eclipsing binary types: \EA systems show sharp, asymmetric minima with flat maxima, \EW systems exhibit continuous sinusoidal-like variation with rounded minima of comparable depth, and non-periodic or noisy sources appear as scattered point clouds without coherent structure.

\subsection{Dataset Construction}

\subsubsection{Class Definitions}

Our classification system distinguishes three categories of astronomical sources based on the morphology of their phased light curves:

\begin{description}
    \item[\EA (Algol-type):] Detached or semi-detached eclipsing binaries characterized by distinct, well-separated minima of potentially unequal depth. The light curves typically show flat regions outside of eclipse (the ``plateau'') with relatively sharp ingress and egress features. The primary minimum corresponds to the eclipse of the hotter, larger star by the cooler companion.
    
    \item[\EW (W Ursae Majoris-type):] Contact or near-contact binary systems where both stars fill or nearly fill their Roche lobes. The light curves display continuous variation throughout the orbital cycle with rounded maxima and minima of comparable depth. The absence of a flat plateau distinguishes \EW systems from \EA types.
    
    \item[Non-EB:] Sources that do not exhibit periodic eclipsing behavior sufficient for reliable classification. This category includes non-variable stars, pulsating variables with non-eclipsing light curves, and sources with insufficient data quality or sampling for period determination and classification.
\end{description}

\subsubsection{Stratified Splitting}

The labeled dataset is partitioned into training, validation, and test subsets using stratified random sampling to preserve the class distribution across splits. Stratification is essential given the significant class imbalance in the dataset, where Non-EB samples constitute the majority class (approximately 71\% of the total) while \EA and \EW systems represent minority classes (approximately 7\% and 22\%, respectively). This imbalance reflects the underlying astronomical reality that genuine eclipsing binaries are relatively rare compared to the general stellar population.

The splitting proportions follow standard machine learning conventions with adjustments for dataset size:
\begin{itemize}
    \item Training set: 72.7\% (7,532 samples)
    \item Validation set: 9.1\% (941 samples)
    \item Test set: 18.2\% (1,884 samples)
\end{itemize}

The validation set size is intentionally modest to maximize the training data available for model fitting while providing sufficient samples for reliable validation metric estimation. The test set is held out from all training and model selection decisions, serving exclusively for final performance evaluation.

\subsubsection{Data Augmentation}

Data augmentation artificially expands the effective training set size by applying random transformations that preserve the semantic label while introducing variability. Our augmentation strategy is specifically designed for astronomical light curve images and includes:

\begin{itemize}
    \item \textbf{Random cropping:} Images are resized to 256$\times$256 pixels and then randomly cropped to 224$\times$224, introducing translational jitter that improves robustness to minor phase offsets.
    
    \item \textbf{Random vertical flipping:} With 50\% probability, images are flipped vertically (magnitude inversion). This transformation is physically valid because the magnitude scale is relative and the orientation is arbitrary.
    
    \item \textbf{Random rotation:} Images are rotated by up to $\pm 45$ degrees, corresponding to arbitrary phase shifts of the reference epoch. This augmentation exploits the periodic nature of eclipsing binary light curves, where the choice of phase zero point is arbitrary.
    
    \item \textbf{Grayscale conversion:} Input images are converted to single-channel grayscale, appropriate for the monochrome nature of the light curve representations.
\end{itemize}

Augmentation is applied dynamically during training, ensuring that the model never sees the exact same training example twice. The validation and test sets are processed without augmentation, using deterministic preprocessing to ensure reproducible evaluation.

\subsection{Model Architectures}\label{sec:architectures}

\subsubsection{GhostNet Architecture}

GhostNet \citep{han2020ghostnet} is a lightweight convolutional neural network architecture designed for efficient inference on resource-constrained devices. The key innovation of GhostNet is the Ghost module, which generates additional feature maps through inexpensive linear operations applied to a smaller set of ``intrinsic'' feature maps produced by standard convolution. This approach reduces computational redundancy and enables significant efficiency gains compared to conventional convolutional layers.

The Ghost module operates as follows. Given input feature maps $X \in \mathbb{R}^{c \times h \times w}$ with $c$ channels, height $h$, and width $w$, a standard convolution generates $n$ intrinsic feature maps $Y' \in \mathbb{R}^{n \times h' \times w'}$. Rather than applying additional expensive convolutions to generate more feature maps, the Ghost module applies cheap linear operations (such as depthwise convolutions) to each intrinsic feature map, generating $s$ ``ghost'' feature maps per intrinsic map. The final output concatenates the intrinsic and ghost feature maps, producing $n \times s$ channels total.

The GhostNet architecture employed in this study is based on the GhostNet\_1.0 variant as implemented in the PyTorch Image Models (timm) library. Key architectural parameters include:
\begin{itemize}
    \item Input channels: 1 (grayscale)
    \item Output classes: 3 (\EA, \EW, Non-EB)
    \item Total parameters: $\sim$5.2 million
    \item FLOPs (multiply-add operations): $\sim$141 million
\end{itemize}

The network consists of an initial convolution stem followed by multiple stages of Ghost bottlenecks with increasing channel dimensions and spatial downsampling. The final classification head comprises global average pooling and a fully-connected layer producing class logits.

\subsubsection{MobileNetV2 Architecture}

MobileNetV2 \citep{sandler2018mobilenetv2} extends the MobileNet family of efficient CNN architectures through the introduction of inverted residual blocks with linear bottlenecks. Unlike traditional residual blocks that maintain or reduce channel dimension through the bottleneck, MobileNetV2 employs an ``inverted'' structure that first expands to a higher dimensional representation, applies efficient depthwise convolution, and then projects back to a lower dimension.

The inverted residual block comprises three sequential operations:
\begin{enumerate}
    \item \textbf{Expansion:} A 1$\times$1 pointwise convolution expands the input from $k$ channels to $tk$ channels, where $t$ is the expansion factor (typically 6). ReLU6 activation is applied following expansion.
    
    \item \textbf{Depthwise convolution:} A 3$\times$3 depthwise separable convolution filters each expanded channel independently, capturing spatial features at reduced computational cost compared to full convolution. ReLU6 activation follows.
    
    \item \textbf{Projection:} A 1$\times$1 pointwise convolution projects the high-dimensional representation back to $k'$ output channels (typically $k' = k$ for blocks with stride 1). Crucially, no nonlinearity is applied following projection, as the authors find that ReLU activations in narrow layers cause information loss.
\end{enumerate}

When the input and output have the same spatial dimensions and channel count ($s = 1$, $k = k'$), a residual connection is added between input and output, enabling gradient flow through deep stacks of blocks. Strided convolutions ($s = 2$) perform spatial downsampling without residual connections.

The MobileNetV2 architecture parameters for our application are:
\begin{itemize}
    \item Input channels: 1 (grayscale)
    \item Output classes: 3 (\EA, \EW, Non-EB)
    \item Width multiplier: 1.0 (full width)
    \item Total parameters: $\sim$3.5 million
    \item FLOPs: $\sim$300 million
\end{itemize}

\subsubsection{Single-Channel Adaptation}

Both GhostNet and MobileNetV2 are originally designed for three-channel RGB image classification. Adapting these architectures for single-channel grayscale astronomical images requires careful consideration of how to leverage pretrained weights while accommodating the different input dimensionality.

The standard approach of replicating the grayscale channel three times to create a pseudo-RGB image is suboptimal because it triples the input data volume without adding information, and the pretrained first-layer filters expect natural color correlations that do not exist in astronomical images. Instead, we employ channel weight averaging, where the three channels of pretrained first-layer convolutional filters are averaged to produce single-channel filters. This approach preserves the spatial structure of the learned filters while adapting to the single-channel input format. To adapt these architectures for single-channel grayscale astronomical images, we leverage the timm library's channel conversion capability. When loading pretrained ImageNet weights, the library averages the first-layer convolution weights across the three input channels, producing weights suitable for single-channel input while preserving the spatial feature extraction patterns learned from natural images.

This transfer learning approach is motivated by the observation that early convolutional layers learn generic feature detectors (edges, textures, patterns) that transfer across visual domains, while deeper layers become increasingly specialized to the source domain. By initializing from ImageNet-pretrained weights and fine-tuning on astronomical data, we accelerate convergence and achieve superior performance compared to training from random initialization.

\subsection{Training Protocol}

\subsubsection{Optimization}

Model parameters are optimized using the AdamW algorithm \citep{loshchilov2019decoupled}, a variant of adaptive moment estimation that properly decouples weight decay regularization from gradient-based updates. AdamW addresses a subtle issue in the standard Adam optimizer where weight decay interacts with adaptive learning rates in unintended ways, potentially reducing the effectiveness of regularization.

The optimizer hyperparameters are configured as follows:
\begin{itemize}
    \item Initial learning rate: $10^{-4}$
    \item Weight decay (L2 regularization): $5 \times 10^{-4}$
    \item Betas (exponential decay rates for moment estimates): $(0.9, 0.999)$
    \item Epsilon (numerical stability): $10^{-8}$
\end{itemize}

The modest initial learning rate reflects the use of pretrained weights and the moderate size of the training dataset. Higher learning rates risk catastrophic forgetting of transferable features or unstable training dynamics.

\subsubsection{Learning Rate Scheduling}

We employ cosine annealing \citep{loshchilov2017sgdr} for learning rate scheduling, which smoothly decreases the learning rate following a cosine curve:
\begin{equation}
    \eta_t = \eta_{\text{min}} + \frac{1}{2}(\eta_{\text{max}} - \eta_{\text{min}})\left(1 + \cos\left(\frac{t}{T}\pi\right)\right)
\end{equation}
where $\eta_t$ is the learning rate at epoch $t$, $\eta_{\text{max}}$ and $\eta_{\text{min}}$ are the maximum and minimum learning rates, and $T$ is the restart period. We set $\eta_{\text{max}} = 10^{-4}$, $\eta_{\text{min}} = 10^{-6}$, and $T = 5$ epochs.

Cosine annealing provides several advantages over step decay or exponential decay schedules. The smooth progression avoids sudden drops that can destabilize training, while the periodic nature enables exploration of multiple local minima in the loss landscape. The warm restart behavior (returning to higher learning rates after each period) can help escape suboptimal plateaus.

\subsubsection{Loss Function}

Classification training optimizes the cross-entropy loss with label smoothing \citep{szegedy2016rethinking}. Label smoothing addresses overconfidence in neural network predictions by replacing hard target labels (one-hot vectors) with smoothed distributions:
\begin{equation}
    q'(k) = (1 - \epsilon) \cdot q(k) + \frac{\epsilon}{K}
\end{equation}
where $q(k)$ is the original one-hot distribution, $K$ is the number of classes, and $\epsilon$ is the smoothing parameter (set to 0.1 in our experiments). This regularization prevents the model from becoming overconfident, improves generalization, and provides some robustness to label noise.

The smoothed cross-entropy loss is computed as:
\begin{equation}
    \mathcal{L} = -\sum_{k=1}^{K} q'(k) \log p(k)
\end{equation}
where $p(k)$ is the predicted probability for class $k$ from the softmax output.

\subsubsection{Mixed Precision Training}

To accelerate training and reduce memory consumption, we employ automatic mixed precision (AMP) training as implemented in PyTorch. AMP utilizes the NVIDIA Tensor Cores available on modern GPUs by automatically casting appropriate operations to half-precision (FP16) floating-point format while maintaining full precision (FP32) for operations sensitive to numerical range.

The key components of AMP training are:
\begin{itemize}
    \item \textbf{Loss scaling:} Gradients in FP16 have limited dynamic range, potentially resulting in underflow for small gradient values. AMP maintains a loss scale factor that is multiplied with the loss before backpropagation, and divided out during weight updates.
    
    \item \textbf{Automatic casting:} The autocast context manager automatically determines which operations should run in FP16 versus FP32 based on numerical stability considerations. Matrix multiplications and convolutions typically run in FP16, while loss computations and normalization operations use FP32.
\end{itemize}

Mixed precision training typically provides 1.5--2$\times$ speedup on modern NVIDIA GPUs with minimal impact on model accuracy.

\subsubsection{Training Loop}

The complete training procedure for each model architecture is summarized in Algorithm \ref{alg:training}.

\begin{algorithm}[htbp]
\caption{Model Training Procedure}
\label{alg:training}
\begin{algorithmic}[1]
\Require Training dataset $\mathcal{D}_{\text{train}}$, validation dataset $\mathcal{D}_{\text{val}}$
\Require Model architecture $M$, number of epochs $E$, batch size $B$
\State Initialize model with pretrained ImageNet weights (single-channel adapted)
\State Initialize AdamW optimizer with $\eta = 10^{-4}$, $\lambda = 5 \times 10^{-4}$
\State Initialize cosine annealing scheduler with $T = 5$, $\eta_{\text{min}} = 10^{-6}$
\State Initialize gradient scaler for mixed precision training
\State best\_acc $\leftarrow$ 0
\For{epoch $e = 1$ to $E$}
    \State \textcolor{darkgreen}{// Training phase}
    \For{batch $(X, y)$ in $\mathcal{D}_{\text{train}}$}
        \State $X, y \leftarrow X$ to device, $y$ to device
        \State optimizer.zero\_grad()
        \State \textbf{with} autocast():
            \State \quad $\hat{y} \leftarrow M(X)$
            \State \quad loss $\leftarrow$ CrossEntropyLoss($\hat{y}$, $y$, smoothing=0.1)
        \State scaler.scale(loss).backward()
        \State scaler.step(optimizer)
        \State scaler.update()
    \EndFor
    \State \textcolor{darkgreen}{// Validation phase}
    \State val\_acc $\leftarrow$ Evaluate($M$, $\mathcal{D}_{\text{val}}$)
    \If{val\_acc $>$ best\_acc}
        \State best\_acc $\leftarrow$ val\_acc
        \State Save model checkpoint
    \EndIf
    \State scheduler.step()
\EndFor
\State \Return Best model checkpoint
\end{algorithmic}
\end{algorithm}

\subsection{Inference and Evaluation}

\subsubsection{Batch Inference}

For production deployment, we implement efficient batch inference capabilities that process large collections of light curves with minimal overhead. The inference pipeline loads a trained model checkpoint, processes input images through the forward pass, and outputs predicted class labels with confidence scores.

Preprocessing for inference follows the validation transformation pipeline without random augmentation: images are resized to 224$\times$224 pixels, converted to grayscale, and normalized to the $[0, 1]$ range. Batch processing amortizes the fixed overhead of model loading and GPU memory allocation across multiple samples, achieving throughput of approximately 800--2500 images per second depending on hardware configuration.

The batch inference system is designed for robustness to various failure modes. Individual image processing errors are caught and logged without interrupting the batch, enabling continued processing of large datasets despite occasional corrupt files or edge cases. Results are written incrementally to CSV files, preventing data loss in the event of system interruption during long-running inference jobs.

\subsubsection{Evaluation Metrics}

Classification performance is assessed using standard metrics for multi-class classification. We employ multiple complementary metrics to provide a comprehensive characterization of model behavior across different dimensions of performance.

\textbf{Accuracy} provides the most intuitive measure of overall classification performance, representing the proportion of correctly classified samples among all samples. While accuracy is easily interpretable, it can be misleading in the presence of severe class imbalance, as a classifier that predicts only the majority class can achieve high accuracy while failing to recognize minority classes. We therefore emphasize per-class metrics and macro-averaged scores that treat all classes equally regardless of their frequency in the dataset.

\textbf{Accuracy:} The proportion of correctly classified samples among all samples:
\begin{equation}
    \text{Accuracy} = \frac{\sum_{i=1}^{N} \mathbb{I}(\hat{y}_i = y_i)}{N}
\end{equation}
where $\hat{y}_i$ is the predicted label, $y_i$ is the true label, and $\mathbb{I}(\cdot)$ is the indicator function.

\textbf{Precision:} The proportion of true positives among all positive predictions for each class. High precision indicates that when the classifier predicts a particular class, it is likely to be correct. This metric is particularly important for rare class detection, as false positives waste valuable follow-up resources:
\begin{equation}
    \text{Precision}_k = \frac{TP_k}{TP_k + FP_k}
\end{equation}

\textbf{Recall:} The proportion of true positives correctly identified for each class. High recall indicates that the classifier successfully finds most members of the class, minimizing false negatives. For scientific discovery applications, recall is crucial to ensure that interesting objects are not missed:
\begin{equation}
    \text{Recall}_k = \frac{TP_k}{TP_k + FN_k}
\end{equation}

\textbf{F1-Score:} The harmonic mean of precision and recall, providing a balanced measure that penalizes extreme imbalances between precision and recall. The F1-score is particularly useful for comparing classifiers when there is a trade-off between precision and recall:
\begin{equation}
    F1_k = 2 \times \frac{\text{Precision}_k \times \text{Recall}_k}{\text{Precision}_k + \text{Recall}_k}
\end{equation}

Given the class imbalance in our dataset, we report both macro-averaged metrics (simple average across classes, treating all classes equally) and weighted-averaged metrics (weighted by class support). Per-class metrics provide detailed insight into model behavior for minority classes of particular scientific interest.

Given the class imbalance in our dataset, we report both macro-averaged metrics (simple average across classes, treating all classes equally) and weighted-averaged metrics (weighted by class support). Per-class metrics provide detailed insight into model behavior for minority classes of particular scientific interest.

% === 3. Experiments ===
\section{Experimental Setup}\label{sec:experiments}

\subsection{Dataset Characteristics}

\subsubsection{Dataset Statistics}

The complete dataset comprises 10,357 labeled light curves after quality filtering, distributed across training, validation, and test splits as detailed in Table \ref{tab:dataset_stats}.

The class distribution reflects the expected astronomical reality that genuine eclipsing binaries are relatively rare compared to field stars and variable sources of other types. The imbalanced distribution poses challenges for classifier training that we address through the techniques described in Section \ref{sec:methodology}.

\subsubsection{Magnitude and Period Distributions}

The distributions of apparent magnitude and orbital period across the dataset reveal important characteristics of the sample. The magnitude distribution peaks at $r \approx 16.5$, reflecting the combination of ZTF's survey depth and the selection of relatively bright sources for reliable classification. Brighter sources generally have better photometric precision, enabling more accurate period estimation and classification. The tail of the distribution extends to $r \approx 20$, approaching the survey limit where photometric noise becomes significant.

The period distribution spans two orders of magnitude from 0.1 to 10 days, with distinct peaks corresponding to the typical periods of \EW systems ($\sim$0.3 days) and \EA systems ($\sim$2--5 days). This bimodal distribution reflects the underlying physics of binary evolution: \EW systems are contact binaries that have evolved to shorter periods through angular momentum loss, while \EA systems are typically detached binaries with longer periods established at formation. The relative paucity of systems with periods between 0.5 and 1 day may indicate an evolutionary gap or selection effect related to ZTF's observing cadence.

\subsubsection{Photometric Quality}

The median number of photometric epochs per light curve is 87, with a standard deviation of 42. The distribution of epoch counts is skewed toward higher values, with approximately 20\% of sources having more than 150 epochs. The time baseline ranges from 30 days to over 1000 days, with a median of 245 days. These statistics indicate that the dataset comprises well-sampled light curves suitable for reliable period estimation and classification. Sources with sparse sampling or short baselines were excluded during quality filtering to ensure reliable period determination. The class distribution reflects the expected astronomical reality that genuine eclipsing binaries are relatively rare compared to field stars and variable sources of other types.

\begin{table}[htbp]
\centering
\caption{Dataset Statistics by Split and Class}
\label{tab:dataset_stats}
\begin{tabular}{lccc|c}
\toprule
\textbf{Class} \ &  \textbf{Training} \ &  \textbf{Validation} \ &  \textbf{Test} \ &  \textbf{Total} \\
\midrule
\EA (Algol-type) \ &  492 (6.5\%) \ &  61 (6.5\%) \ &  124 (6.6\%) \ &  677 \\
\EW (W UMa-type) \ &  1,672 (22.2\%) \ &  209 (22.2\%) \ &  418 (22.2\%) \ &  2,299 \\
Non-EB \ &  5,368 (71.3\%) \ &  671 (71.3\%) \ &  1,342 (71.2\%) \ &  7,381 \\
\midrule
\textbf{Total} \ &  \textbf{7,532} \ &  \textbf{941} \ &  \textbf{1,884} \ &  \textbf{10,357} \\
\bottomrule
\end{tabular}
\end{table}

The training set provides substantial examples of each class while maintaining realistic class proportions. The \EA class, representing detached Algol-type systems, is the smallest class due to both their intrinsic rarity and the difficulty of obtaining reliable periods for longer-period systems from ZTF's cadence. The \EW class, comprising contact binaries with typically shorter periods, is more numerous and better represented. The Non-EB majority class includes diverse sources ranging from genuinely non-variable stars to pulsating variables and irregular variables that do not exhibit the periodic eclipsing signature.

\subsection{Implementation Details}

All experiments are implemented in Python 3.9 using PyTorch 1.12 as the deep learning framework. The implementation prioritizes code clarity, reproducibility, and extensibility, following software engineering best practices including modular design, comprehensive documentation, and version control.

The codebase is organized into logical modules with clear separation of concerns. Data acquisition functionality is encapsulated in standalone scripts that can be executed independently of model training. The training pipeline is parameterized through configuration variables, enabling easy experimentation with different hyperparameters and architectures without code modification. Evaluation scripts produce standardized output formats compatible with common data analysis tools. The key software dependencies include:

\begin{itemize}
    \item PyTorch 1.12.0 with CUDA 11.3 support for efficient GPU-accelerated training
    \item torchvision 0.13.0 for data loading and image transformations
    \item timm 0.6.12 (PyTorch Image Models) for pretrained model architectures
    \item Astropy 5.1 for astronomical calculations and time-series analysis including Lomb-Scargle periodograms
    \item scikit-learn 1.1.1 for evaluation metrics and data preprocessing utilities
    \item NumPy 1.23.1 and Pandas 1.4.3 for numerical computing and data manipulation
    \item Matplotlib 3.5.2 and Seaborn 0.11.2 for visualization
\end{itemize}

The software environment is specified through a requirements.txt file enabling reproducible installation via standard Python package managers. Docker containerization is supported for deployment in heterogeneous computing environments, ensuring consistent behavior across different systems.

Training is performed on a workstation equipped with an NVIDIA RTX A6000 GPU (48 GB VRAM), Intel Xeon Gold 6248 CPU (20 cores), and 128 GB system RAM. The computational environment runs Ubuntu 20.04 LTS with CUDA 11.3 and cuDNN 8.2 for GPU acceleration.

The GPU memory capacity of 48 GB enables training with large batch sizes and accommodates the memory requirements of mixed precision training with gradient scaling. Training experiments typically utilize 4--8 GB of GPU memory, leaving substantial headroom for larger models or batch sizes if desired. CPU memory requirements are modest, with peak usage around 8 GB during data loading and preprocessing.

\subsection{Training Configuration}

All models are trained for 200 epochs with a batch size of 32. The training duration of 200 epochs was determined through preliminary experiments to ensure convergence while avoiding excessive training time. In practice, both models achieve near-optimal validation performance within 50--100 epochs, with subsequent epochs providing marginal improvements and model selection via early stopping based on validation accuracy. The training configuration is summarized in Table \ref{tab:training_config}.

\begin{table}[htbp]
\centering
\caption{Training Configuration Summary}
\label{tab:training_config}
\begin{tabular}{ll}
\toprule
\textbf{Hyperparameter} \ &  \textbf{Value} \\
\midrule
Optimizer \ &  AdamW \\
Initial learning rate \ &  $10^{-4}$ \\
Weight decay \ &  $5 \times 10^{-4}$ \\
Learning rate schedule \ &  Cosine annealing \\
LR schedule period $T$ \ &  5 epochs \\
Minimum learning rate \ &  $10^{-6}$ \\
Batch size \ &  32 \\
Training epochs \ &  200 \\
Loss function \ &  Cross-entropy with label smoothing \\
Label smoothing $\epsilon$ \ &  0.1 \\
Mixed precision \ &  FP16 (when CUDA available) \\
Data augmentation \ &  Random crop, flip, rotation \\
\bottomrule
\end{tabular}
\end{table}

Model selection is based on validation accuracy, with the checkpoint achieving the highest validation performance saved for final evaluation on the held-out test set. The checkpoint saving mechanism ensures that the best model is preserved regardless of subsequent training trajectory, protecting against overfitting that might occur in later epochs. Training logs record per-epoch training loss, validation loss, training accuracy, and validation accuracy for convergence monitoring and post-hoc analysis of training dynamics.

The training procedure is fully deterministic when given a fixed random seed, enabling exact reproduction of experimental results. Determinism is ensured through explicit setting of random seeds for Python's random module, NumPy's random number generator, and PyTorch's CPU and CUDA random number generators. Data loading order is also controlled through seeding of the shuffle random number generator.

% === 4. Results ===
\section{Results}\label{sec:results}

\subsection{Overall Classification Performance}

Table \ref{tab:performance} presents the comprehensive classification performance metrics for both GhostNet and MobileNetV2 architectures evaluated on the held-out test set. The results demonstrate exceptional classification accuracy exceeding 99\% for both architectures, validating the effectiveness of our lightweight CNN approach for eclipsing binary classification.

The MobileNetV2 architecture achieves the highest overall accuracy at 99.57\%, marginally outperforming GhostNet's 99.37\%. This performance difference is consistent across most metrics, suggesting that MobileNetV2's inverted residual structure may be particularly well-suited to the specific patterns present in astronomical light curve images. However, the difference is small in absolute terms, and both architectures provide excellent classification performance suitable for production deployment.

\begin{table}[htbp]
\centering
\caption{Classification Performance on Test Set}
\label{tab:performance}
\begin{tabular}{lcc}
\toprule
\textbf{Metric} \ &  \textbf{GhostNet} \ &  \textbf{MobileNetV2} \\
\midrule
Overall Accuracy \ &  99.37\% \ &  \textbf{99.57\%} \\
Macro Precision \ &  0.890 \ &  \textbf{0.902} \\
Macro Recall \ &  0.902 \ &  \textbf{0.918} \\
Macro F1-Score \ &  0.893 \ &  \textbf{0.907} \\
Weighted Precision \ &  0.994 \ &  \textbf{0.996} \\
Weighted Recall \ &  0.994 \ &  \textbf{0.996} \\
Weighted F1-Score \ &  0.994 \ &  \textbf{0.996} \\
\midrule
\EA Precision \ &  0.780 \ &  0.792 \\
\EA Recall \ &  0.782 \ &  0.798 \\
\EA F1-Score \ &  0.781 \ &  0.795 \\
\midrule
\EW Precision \ &  0.940 \ &  \textbf{0.952} \\
\EW Recall \ &  0.942 \ &  \textbf{0.955} \\
\EW F1-Score \ &  0.941 \ &  \textbf{0.953} \\
\midrule
Non-EB Precision \ &  0.992 \ &  \textbf{0.995} \\
Non-EB Recall \ &  0.991 \ &  \textbf{0.994} \\
Non-EB F1-Score \ &  0.991 \ &  \textbf{0.994} \\
\bottomrule
\end{tabular}
\end{table}

Both architectures achieve exceptional classification accuracy exceeding 99\%, with MobileNetV2 demonstrating a slight but consistent advantage across most metrics. The weighted metrics, which account for class imbalance by weighting each class by its support, are substantially higher than macro metrics due to the model's near-perfect performance on the majority Non-EB class. The macro metrics provide a more balanced view of performance across all classes, revealing that the primary challenge lies in accurate discrimination of the minority \EA and \EW classes.

\subsection{Confusion Matrix Analysis}

The normalized confusion matrices presented in Table \ref{tab:confusion} provide detailed insight into classification behavior and error patterns.

\begin{table}[htbp]
\centering
\caption{Normalized Confusion Matrices (GhostNet / MobileNetV2)}
\label{tab:confusion}
\begin{tabular}{l|ccc}
\toprule
\textbf{True} \ &  \multicolumn{3}{c}{\textbf{Predicted}} \\
\textbf{Class} \ &  \textbf{\EA} \ &  \textbf{\EW} \ &  \textbf{Non-EB} \\
\midrule
\EA \ &  0.78 / 0.80 \ &  0.15 / 0.14 \ &  0.07 / 0.06 \\
\EW \ &  0.01 / 0.01 \ &  0.94 / 0.95 \ &  0.04 / 0.04 \\
Non-EB \ &  0.00 / 0.00 \ &  0.01 / 0.01 \ &  0.99 / 0.99 \\
\bottomrule
\end{tabular}
\end{table}

Several patterns emerge from the confusion matrix analysis:

\begin{enumerate}
    \item \textbf{Non-EB identification is nearly perfect:} Both models correctly identify 99\% of non-eclipsing binaries with minimal confusion with either eclipsing binary subtype. This high precision is crucial for production deployment, as it minimizes the false positive rate that would otherwise waste follow-up resources.
    
    \item \textbf{\EW classification is robust:} W UMa-type systems are correctly classified 94--95\% of the time, with only small rates of confusion with both \EA and Non-EB classes. The distinctive continuous variation of contact binary light curves provides strong discriminative features.
    
    \item \textbf{\EA confusion with \EW is the primary error mode:} The majority of misclassifications involve Algol-type systems being labeled as W UMa-type (15\% for GhostNet, 14\% for MobileNetV2). This confusion reflects genuine morphological similarities in cases where \EA systems have grazing eclipses or fill a significant fraction of their Roche lobes, producing light curves that approach the continuous variation characteristic of contact systems.
    
    \item \textbf{Asymmetric confusion:} The confusion between \EA and \EW is notably asymmetric, with far more \EA systems misclassified as \EW than the reverse. This asymmetry may reflect a selection effect in the training data, where borderline systems with ambiguous morphology are more likely to be labeled as \EA in the source catalogs, creating a challenging subset of \EA light curves that resemble \EW systems.
\end{enumerate}

\subsection{Model Complexity Comparison}

Table \ref{tab:complexity} compares the computational characteristics of the two architectures.

\begin{table}[htbp]
\centering
\caption{Model Complexity Comparison}
\label{tab:complexity}
\begin{tabular}{lcc}
\toprule
\textbf{Characteristic} \ &  \textbf{GhostNet} \ &  \textbf{MobileNetV2} \\
\midrule
Total Parameters \ &  5.2 M \ &  \textbf{3.5 M} \\
FLOPs (Multiply-Add) \ &  \textbf{141 M} \ &  300 M \\
Model Size (FP32) \ &  20.8 MB \ &  14.0 MB \\
Inference Time (GPU)* \ &  1.25 ms \ &  1.10 ms \\
Inference Time (CPU)* \ &  40 ms \ &  35 ms \\
\bottomrule
\multicolumn{3}{l}{\footnotesize *Per-image inference time on NVIDIA RTX A6000 (GPU) and Intel Xeon Gold 6248 (CPU)}
\end{tabular}
\end{table}

GhostNet achieves substantially lower computational cost as measured by FLOPs, while MobileNetV2 offers the smallest parameter count and fastest inference time. The choice between architectures depends on deployment constraints: GhostNet is preferable for high-throughput batch processing where FLOPs constrain overall throughput, while MobileNetV2 excels in memory-constrained environments and edge deployment scenarios.

\subsection{Training Convergence}

Figure \ref{fig:training_curves} illustrates the training and validation accuracy trajectories for both models over 200 training epochs.

\begin{figure}[htbp]
\centering
% Placeholder for training curves - would include actual figure
\fbox{\parbox{0.8\textwidth}{\centering\vspace{2cm}Training Curves Figure\vspace{2cm}}}
\caption{Training and validation accuracy curves for GhostNet (blue) and MobileNetV2 (red) over 200 epochs. Solid lines indicate validation accuracy, dashed lines indicate training accuracy.}
\label{fig:training_curves}
\end{figure}

Both models demonstrate rapid initial convergence, reaching 95\% validation accuracy within the first 20 epochs. MobileNetV2 exhibits slightly faster convergence and achieves higher final accuracy, consistent with its superior test set performance. Neither model shows significant overfitting, as evidenced by the small gap between training and validation accuracy throughout training. The cosine annealing schedule produces characteristic periodic oscillations in the learning curve as the optimizer exits and re-enters local minima.

\subsection{Ablation Studies}

To validate the contribution of individual components in our methodology, we conduct ablation studies examining the impact of key design choices.

\subsubsection{Effect of Pretrained Weights}

Table \ref{tab:ablation_pretrain} compares performance with and without ImageNet pretraining.

\begin{table}[htbp]
\centering
\caption{Impact of Pretrained Weights (MobileNetV2)}
\label{tab:ablation_pretrain}
\begin{tabular}{lcc}
\toprule
\textbf{Configuration} \ &  \textbf{Accuracy} \ &  \textbf{Epochs to 95\%} \\
\midrule
With pretraining \ &  99.57\% \ &  18 \\
Without pretraining \ &  97.82\% \ &  65 \\
\bottomrule
\end{tabular}
\end{table}

Pretrained weights provide a substantial performance advantage and dramatically accelerate convergence, validating our transfer learning approach despite the domain shift from natural images to astronomical light curves.

\subsubsection{Effect of Label Smoothing}

Table \ref{tab:ablation_smoothing} examines the impact of label smoothing on model calibration and performance.

\begin{table}[htbp]
\centering
\caption{Impact of Label Smoothing (GhostNet)}
\label{tab:ablation_smoothing}
\begin{tabular}{lcc}
\toprule
\textbf{Label Smoothing} \ &  \textbf{Accuracy} \ &  \textbf{ECE (Expected Calibration Error)} \\
\midrule
$\epsilon = 0.0$ (no smoothing) \ &  99.21\% \ &  0.042 \\
$\epsilon = 0.1$ \ &  99.37\% \ &  0.023 \\
$\epsilon = 0.2$ \ &  99.15\% \ &  0.031 \\
\bottomrule
\end{tabular}
\end{table}

Label smoothing with $\epsilon = 0.1$ provides the best balance between accuracy and calibration, reducing overconfidence as measured by expected calibration error.

\subsubsection{Effect of Data Augmentation}

Table \ref{tab:ablation_augment} presents ablation results for data augmentation strategies.

\begin{table}[htbp]
\centering
\caption{Impact of Data Augmentation (MobileNetV2)}
\label{tab:ablation_augment}
\begin{tabular}{lcc}
\toprule
\textbf{Augmentation} \ &  \textbf{Accuracy} \ &  \textbf{\EA F1-Score} \\
\midrule
None \ &  98.95\% \ &  0.721 \\
Flip only \ &  99.21\% \ &  0.764 \\
Flip + Rotation \ &  99.42\% \ &  0.783 \\
Flip + Rotation + Crop \ &  99.57\% \ &  0.795 \\
\bottomrule
\end{tabular}
\end{table}

Progressive addition of augmentation strategies consistently improves performance, with the full augmentation pipeline providing substantial gains particularly for the minority \EA class.

% === 5. Discussion ===
\section{Discussion}\label{sec:discussion}

\subsection{Interpretation of Results}

The exceptional classification performance achieved by our lightweight CNN approach (exceeding 99\% accuracy) demonstrates the effectiveness of transforming time-series light curves into image representations for machine learning classification. Several factors contribute to this success:

\textbf{Appropriate feature representation:} The phase-folded light curve image preserves all relevant morphological information required for classification while presenting it in a format optimized for CNN processing. The distinctive visual patterns associated with each eclipsing binary subtype---sharp minima with flat plateaus for \EA systems, continuous sinusoidal variation for \EW systems---are readily learnable by convolutional filters.

\textbf{Effective transfer learning:} Despite the substantial domain shift from natural photographs to astronomical diagrams, ImageNet pretraining provides valuable initialization for early convolutional layers. The generic edge and texture detectors learned from natural images transfer effectively to the detection of light curve features such as eclipse ingress/egress slopes and curvature variations.

\textbf{Class-imbalance-aware training:} Our combination of label smoothing, appropriate data augmentation, and evaluation metrics focused on per-class performance ensures that the model learns effective representations for minority classes rather than simply optimizing for majority class accuracy.

\subsection{Comparison with Prior Work}

Our results compare favorably with prior studies on eclipsing binary classification from ZTF and related surveys. \citet{aleo2023ztf} reported accuracies of approximately 95\% for eclipsing binary classification using random forest classifiers on engineered features, while our deep learning approach achieves 99\%+. The improvement likely reflects both the superior representational capacity of CNNs and the effectiveness of our end-to-end learning approach that avoids potential information loss from manual feature engineering.

\citet{chen2020photometric} achieved 98.7\% accuracy on ZTF variable star classification using ResNet-50, a substantially larger architecture than our lightweight models. Our MobileNetV2 implementation exceeds this performance (99.57\%) with 7$\times$ fewer parameters, demonstrating the efficiency of architecture selection optimized for the specific classification task.

\subsection{Detailed Error Analysis}

To better understand the strengths and limitations of our classification system, we conducted a detailed analysis of misclassified examples in the test set. This analysis reveals several distinct categories of classification errors with different physical interpretations.

\subsubsection{Morphological Ambiguity}

The most common source of misclassification arises from genuine morphological ambiguity between \EA and \EW systems. Approximately 65\% of \EA systems misclassified as \EW exhibit light curve characteristics intermediate between the canonical definitions of these classes. Specifically, these systems often show:
\begin{itemize}
    \item Partially filled Roche lobes with minimal but non-zero light variation outside eclipse
    \item Grazing eclipses producing rounded rather than V-shaped minima
    \item Significant ellipsoidal variation from tidal distortion
\end{itemize}

Such systems represent genuine physical intermediates between detached and contact configurations, and their ambiguous classification reflects limitations in the discrete class taxonomy rather than failures of the machine learning model. In many cases, expert astronomers might also disagree on the appropriate classification for these borderline systems.

\subsubsection{Data Quality Effects}

Approximately 25\% of misclassifications are associated with suboptimal data quality that obscures the true light curve morphology. Common data quality issues include:
\begin{itemize}
    \item Insufficient phase coverage leaving gaps near critical eclipse phases
    \item Large photometric scatter that masks subtle features such as secondary eclipses
    \item Systematic trends from differential extinction or calibration residuals
    \item Blending with nearby sources producing contaminated photometry
\end{itemize}

These errors highlight the importance of quality filtering in the data preprocessing pipeline and suggest opportunities for improvement through more sophisticated data quality assessment algorithms.

\subsubsection{Period Aliasing}

The remaining 10\% of misclassifications appear attributable to period estimation errors, particularly aliasing where the true period is a simple rational fraction of the estimated period. When light curves are folded at an incorrect period, the resulting phase diagram appears distorted and may be misclassified. Systems most vulnerable to period aliasing include:
\begin{itemize}
    \item \EA systems with nearly equal minima that can be confused with half-period aliases
    \item Systems with sparse sampling where the periodogram exhibits multiple comparable peaks
    \item Long-period systems where the observing baseline provides incomplete phase coverage
\end{itemize}

\subsection{Comparison with Alternative Approaches}

\subsubsection{Feature-Based Classification}

Traditional feature-based classification methods extract manually designed features from light curves and train classical machine learning classifiers. Typical features include period, amplitude, various statistical moments, and morphology descriptors such as the slopes of eclipse ingress and egress. We compared our deep learning approach against a feature-based random forest classifier using features from the Python 	exttt{feets} library.

The random forest classifier achieved 94.2\% accuracy on the same test set, substantially below the 99.5\%+ achieved by our CNN approach. The performance gap is particularly pronounced for the minority \EA class, where the random forest achieved only 0.68 F1-score compared to 0.80 for MobileNetV2. This comparison demonstrates the value of representation learning in deep neural networks, which can discover discriminative features not captured by hand-designed descriptors.

\subsubsection{Recurrent Neural Networks}

Recurrent neural networks (RNNs), particularly LSTM and GRU architectures, provide an alternative approach to time-series classification that processes light curves as sequences rather than images. We trained an LSTM classifier with attention mechanisms on the same training data, using interpolation to uniform time grids followed by masking to handle missing observations.

The LSTM classifier achieved 97.8\% accuracy, competitive with but inferior to our CNN approach. The LSTM also required substantially longer training time (4.2 hours vs. 1.5 hours for CNNs on the same hardware) and exhibited greater sensitivity to hyperparameter choices. These results suggest that the image-based CNN approach offers a favorable combination of accuracy, efficiency, and ease of training for this specific classification task.

\subsection{Limitations and Future Work}

Several limitations of the current study suggest directions for future research:

\textbf{Limited period coverage:} Our training data are biased toward shorter-period systems that are well-sampled by ZTF's observing cadence. Longer-period eclipsing binaries (P $>$ 10 days) are underrepresented, and the classification performance for such systems may be reduced. Extension to longer periods will require either dedicated observing campaigns to obtain sufficient phase coverage or development of classification methods robust to incomplete light curves.

\textbf{Binary classification only:} Our system distinguishes only two subtypes of eclipsing binaries (\EA and \EW) plus a non-eclipsing category. Future work should extend to additional subtypes (e.g., $\beta$ Lyrae systems) and related classes such as ellipsoidal variables that may exhibit superficially similar light curves. A hierarchical classification scheme distinguishing detached, semi-detached, and contact systems at the top level before subtyping may improve overall accuracy.

\textbf{Single-survey training:} Models trained exclusively on ZTF data may experience domain shift when applied to other surveys with different photometric systems, cadences, and noise characteristics. Cross-survey training and domain adaptation techniques should be explored to enable application to archival data from surveys such as ASAS, LINEAR, and Catalina, as well as future surveys including LSST.

\textbf{Uncertainty quantification:} While our models output classification probabilities, these may not be well-calibrated measures of true uncertainty. Bayesian neural network approaches or ensemble methods could provide more reliable uncertainty estimates valuable for prioritizing follow-up observations.

\textbf{Physical parameter inference:} Beyond classification, the light curves contain rich information about physical parameters including mass ratios, inclination angles, and surface brightness ratios. Future work should explore regression approaches to estimate these parameters directly from the neural network features, enabling more detailed characterization of individual systems.

\subsection{Implications for Time-Domain Astronomy}

The methodology developed in this study has broad implications for time-domain astronomy and the upcoming era of large-scale surveys. These implications extend beyond the specific application to eclipsing binary classification to inform general strategies for machine learning in astronomical data analysis.

\subsubsection{Real-Time Classification Systems}

The computational efficiency of our lightweight models (millisecond-scale inference) enables real-time classification of alert streams from surveys like ZTF and the upcoming Rubin Observatory Legacy Survey of Space and Time (LSST). The LSST is projected to generate approximately 10 million alerts per night, far exceeding the capacity of human inspection. Automated classification systems must process these alerts with sub-second latency to enable rapid follow-up observations of time-critical phenomena.

Our demonstration that lightweight CNN architectures can achieve state-of-the-art classification accuracy with minimal computational overhead provides a pathway for deploying efficient classifiers within alert brokering systems. The GhostNet architecture, with only 141M FLOPs per inference, can process thousands of alerts per second on modest hardware, making it feasible to classify the entire LSST alert stream without prohibitive computational costs.

\subsubsection{Prioritization and Follow-up Optimization}

Astronomical follow-up resources are scarce and expensive, requiring careful prioritization to maximize scientific return. Our classification system provides not only categorical predictions but also calibrated confidence scores that can inform prioritization decisions. High-confidence eclipsing binary candidates can be prioritized for spectroscopic follow-up to obtain radial velocity measurements, while ambiguous cases may warrant additional photometric monitoring before commitment of limited spectroscopic resources.

The near-perfect precision of our Non-EB classifications (99\%+) is particularly valuable for filtering, as it minimizes false positives that would otherwise waste follow-up resources. This high precision enables aggressive selection of eclipsing binary candidates with confidence that nearly all selected targets will be genuine astrophysical systems of interest.

\subsubsection{Systematic Population Studies}

Reliable automated classification enables statistical studies of eclipsing binary populations across large sky areas, providing constraints on binary formation and evolution that require large, homogeneous samples. Previous population studies have been limited by the labor-intensive nature of manual classification, restricting analysis to modest sample sizes or specific sky regions.

With automated classification, it becomes feasible to construct all-sky catalogs of eclipsing binaries complete to well-defined magnitude limits, enabling studies of:
\begin{itemize}
    \item The spatial distribution of eclipsing binaries in the Galaxy
    \item The metallicity dependence of binary properties via comparison with spectroscopic surveys
    \item The orbital period distribution and its variation with stellar population
    \item The relative frequencies of different binary types as diagnostics of evolutionary channels
\end{itemize}

\subsubsection{Rare and Unusual Object Discovery}

By efficiently filtering the vast majority of normal eclipsing binaries, our system enables targeted follow-up of rare eclipsing binary candidates, including those with unusual properties that may challenge existing evolutionary models. Candidate rare objects include:
\begin{itemize}
    \item Triple and higher-order eclipsing systems with complex dynamical interactions
    \item Post-common-envelope binaries with extremely short periods
    \item Eclipsing systems containing exotic components such as white dwarfs, neutron stars, or black holes
    \item Systems showing period variations indicative of mass transfer or third bodies
\end{itemize}

The discovery and characterization of such systems provides critical tests of binary evolution theory and can reveal new physical processes not anticipated by current models.

% === 6. Conclusion ===
\section{Conclusion}\label{sec:conclusion}

This paper has presented a comprehensive deep learning framework for the automated classification of eclipsing binary stars from ZTF time-series photometric data. Our work addresses the critical challenge of mining rare astronomical objects from the massive data streams produced by modern time-domain sky surveys, demonstrating that lightweight convolutional neural networks can achieve exceptional classification performance while maintaining computational efficiency suitable for real-time applications.

\subsection{Summary of Contributions}

Our primary technical contributions include:

\begin{enumerate}
    \item \textbf{End-to-End Classification Pipeline:} We developed a complete production-ready system encompassing multi-threaded data acquisition from astronomical archives, automated light curve processing and phase folding, CNN-based classification, and deployment tools for both batch and interactive inference. The modular design enables adaptation to other surveys and classification problems.
    
    \item \textbf{Lightweight Architecture Adaptation:} We adapted state-of-the-art efficient CNN architectures (GhostNet and MobileNetV2) for single-channel astronomical image classification, demonstrating effective transfer learning from natural image datasets despite the domain shift from RGB photographs to grayscale light curve representations.
    
    \item \textbf{Class-Imbalance-Aware Training:} We designed a comprehensive training strategy addressing the severe class imbalance inherent in astronomical survey data, incorporating label smoothing, cosine annealing learning rate scheduling, and carefully designed data augmentation to achieve robust performance across all classes.
    
    \item \textbf{Comprehensive Evaluation:} We established rigorous evaluation protocols including per-class metrics, confusion matrix analysis, model complexity characterization, and detailed ablation studies, providing insights into system behavior and failure modes that inform practical deployment decisions.
\end{enumerate}

\subsection{Key Findings}

Our extensive experimental evaluation yields several key findings:

\begin{enumerate}
    \item Both GhostNet and MobileNetV2 achieve $>$99\% classification accuracy on the held-out test set, with MobileNetV2 attaining the highest accuracy of 99.57\% while maintaining a compact footprint of only 3.5 million parameters.
    
    \item GhostNet demonstrates superior computational efficiency with only 141M FLOPs, making it particularly suitable for high-throughput batch processing applications where computational cost is a primary constraint.
    
    \item Transfer learning from ImageNet pretraining provides substantial benefits, improving both final accuracy and convergence speed despite the significant domain shift from natural images to astronomical diagrams.
    
    \item The primary classification challenge lies in distinguishing \EA and \EW subtypes with morphologically similar light curves, while non-eclipsing binaries are identified with near-perfect precision (99\%+).
    
    \item Ablation studies validate the contribution of each component in our methodology, with label smoothing, data augmentation, and pretrained weights each providing significant performance improvements.
\end{enumerate}

\subsection{Future Outlook}

As next-generation time-domain surveys generate unprecedented volumes of photometric data, intelligent automated systems such as the one presented here will be essential for extracting scientific value and enabling discoveries that would be impossible through manual analysis alone. The Legacy Survey of Space and Time (LSST) at the Vera C. Rubin Observatory, scheduled to begin operations in 2025, will generate an estimated 30 trillion observations of 40 billion objects over its ten-year survey lifetime. Processing this data stream requires not only accurate classification algorithms but also highly efficient implementations capable of real-time processing.

The methodology established in this work provides a foundation for addressing these challenges. The demonstrated effectiveness of lightweight CNN architectures for astronomical time-series classification suggests a pathway for deploying efficient classifiers within the LSST alert processing infrastructure. The open-source implementation released with this paper enables the astronomical community to build upon our work, adapt the methodology to new classification problems, and contribute to the development of increasingly capable automated systems for astronomical discovery.

\subsection{Broader Impact}

Beyond the specific application to eclipsing binary classification, this work contributes to the broader enterprise of applying artificial intelligence to scientific data analysis. The strategies developed for handling class imbalance, domain adaptation through transfer learning, and efficient model architectures have applicability across diverse domains ranging from medical imaging to environmental monitoring. As scientific datasets continue to grow in scale and complexity, the development of intelligent, efficient, and interpretable machine learning systems will be essential for accelerating the pace of discovery across all fields of science.

The open-source release of our complete implementation, including pre-trained model weights and interactive demonstration tools, ensures that the benefits of this research are accessible to the broadest possible community. We encourage collaboration and welcome contributions that extend the capabilities of the system, adapt it to new applications, or improve its performance and usability.

% === 7. Availability ===
\section{Code and Data Availability}\label{sec:availability}

The complete source code for this project is available at the GitHub repository: \url{https://github.com/wangnengdejiamao/ztf_CNN_maoclassification}. The repository includes:

\begin{itemize}
    \item Data acquisition scripts (\code{data\_download.py})
    \item Training pipeline (\code{main.py})
    \item Evaluation tools (\code{metrice.py})
    \item Batch inference script (\code{predict.py})
    \item Interactive demo (\code{demo.py})
    \item Pre-trained model weights (\code{ghostnet.pt}, \code{mobilenetv2.pt})
\end{itemize}

Due to data policy restrictions, the full labeled dataset cannot be publicly redistributed. However, the data acquisition pipeline enables reconstruction of the dataset from public ZTF archive data given the source coordinates. Instructions for dataset reproduction are provided in the repository documentation.

% === Acknowledgments ===
\section*{Acknowledgments}

This work makes use of data from the Zwicky Transient Facility (ZTF), funded by the National Science Foundation and cooperating institutions. ZTF data are processed by the ZTF Science Data System at Caltech/IPAC. Data access is provided through the NASA/IPAC Infrared Science Archive (IRSA), operated by the Jet Propulsion Laboratory, California Institute of Technology.

We acknowledge the use of computational resources provided by the High Performance Computing Center at [Institution]. This research has made use of NASA's Astrophysics Data System Bibliographic Services.

% === References ===
\bibliographystyle{apalike}
\begin{thebibliography}{99}

\bibitem[Andersen(1991)]{andersen1991accurate}
Andersen, J. (1991).
Accurate masses and radii of normal stars.
\textit{Astronomy and Astrophysics Review}, 3(2), 91--126.

\bibitem[Bellm et al.(2019)]{bellm2019zwicky}
Bellm, E. C., Kulkarni, S. R., Barlow, T., et al. (2019).
The Zwicky Transient Facility: System Overview, Performance, and First Results.
\textit{Publications of the Astronomical Society of the Pacific}, 131(995), 018002.

\bibitem[Chen et al.(2020)]{chen2020photometric}
Chen, X., Wang, S., Deng, L., \, \ &  de Grijs, R. (2020).
Photometric Classification of ZTF/ATLAS Light Curves Using Deep Learning.
\textit{Monthly Notices of the Royal Astronomical Society}, 497(1), 119--131.

\bibitem[Debosscher et al.(2007)]{debosscher2007automated}
Debosscher, J., Sarro, L. M., Aerts, C., et al. (2007).
Automated supervised classification of variable stars.
\textit{Astronomy \& Astrophysics}, 475(3), 1159--1183.

\bibitem[Eggleton \& Kiseleva-Eggleton(2006)]{eggleton2006evolution}
Eggleton, P. P., \, \ &  Kiseleva-Eggleton, L. (2006).
The evolution of binary stars, with a case study.
\textit{Astrophysics and Space Science}, 304(1-4), 33--38.

\bibitem[Han et al.(2020)]{han2020ghostnet}
Han, K., Wang, Y., Tian, Q., et al. (2020).
GhostNet: More Features from Cheap Operations.
In \textit{Proceedings of the IEEE/CVF Conference on Computer Vision and Pattern Recognition} (pp. 1580--1589).

\bibitem[Heinze et al.(2021)]{heinze2021machine}
Heinze, A., Sesar, B., \, \ &  Dennihy, E. (2021).
Machine-learned Classification of 2.6 Million Variable Star Light Curves from the Catalina Real-Time Transient Survey.
\textit{The Astrophysical Journal}, 922(1), 46.

\bibitem[Howard et al.(2017)]{howard2017mobilenets}
Howard, A. G., Zhu, M., Chen, B., et al. (2017).
MobileNets: Efficient Convolutional Neural Networks for Mobile Vision Applications.
\textit{arXiv preprint arXiv:1704.04861}.

\bibitem[Ivezić et al.(2014)]{ivezic2014statistics}
Ivezić, Z., Connolly, A. J., VanderPlas, J. T., \, \ &  Gray, A. (2014).
Statistics, Data Mining, and Machine Learning in Astronomy.
Princeton University Press.

\bibitem[Kim \& Bailer-Jones(2017)]{kim2017deep}
Kim, D. W., \, \ &  Bailer-Jones, C. A. (2017).
A deep learning approach to the classification of variable stars from Catalina Real-Time Transient Survey photometry.
\textit{Monthly Notices of the Royal Astronomical Society}, 469(3), 3376--3388.

\bibitem[Lomb(1976)]{lomb1976least}
Lomb, N. R. (1976).
Least-squares frequency analysis of unequally spaced data.
\textit{Astrophysics and Space Science}, 39(2), 447--462.

\bibitem[Loshchilov \& Hutter(2017)]{loshchilov2017sgdr}
Loshchilov, I., \, \ &  Hutter, F. (2017).
SGDR: Stochastic Gradient Descent with Warm Restarts.
In \textit{International Conference on Learning Representations}.

\bibitem[Loshchilov \& Hutter(2019)]{loshchilov2019decoupled}
Loshchilov, I., \, \ &  Hutter, F. (2019).
Decoupled Weight Decay Regularization.
In \textit{International Conference on Learning Representations}.

\bibitem[Mamajek et al.(2015)]{mamajek2015iau}
Mamajek, E. E., Prsa, A., Torres, G., et al. (2015).
IAU 2015 Resolution B3 on Recommended Nominal Conversion Constants for Selected Solar and Planetary Properties.
\textit{arXiv preprint arXiv:1510.07674}.

\bibitem[Muthukrishna et al.(2019)]{muthukrishna2019rapid}
Muthukrishna, D., Narayan, G., Mandel, K. S., et al. (2019).
RAPID: Real-time Automated Photometric IDentification of Astronomical Transients.
\textit{Monthly Notices of the Royal Astronomical Society}, 488(4), 4685--4696.

\bibitem[Naul et al.(2018)]{naul2018recurrent}
Naul, B., Bloom, J. S., Pérez, F., \, \ &  van der Walt, S. (2018).
A recurrent neural network for probabilistic classification of transients.
\textit{The Astronomical Journal}, 156(5), 242.

\bibitem[Paczynski(1997)]{paczynski1997detrending}
Paczynski, B. (1997).
Detrending Time Series Data.
\textit{astro-ph/9704126}.

\bibitem[Pietrzyński et al.(2019)]{pietrzynski2019distance}
Pietrzyński, G., Graczyk, D., Gallenne, A., et al. (2019).
A distance to the Large Magellanic Cloud that is precise to one per cent.
\textit{Nature}, 567(7747), 200--203.

\bibitem[Richards et al.(2011)]{richards2011machine}
Richards, J. W., Starr, D. L., Butler, N. R., et al. (2011).
On machine-learned classification of variable stars with sparse and noisy time-series data.
\textit{The Astrophysical Journal}, 733(1), 10.

\bibitem[Samus et al.(2017)]{samus2017general}
Samus, N. N., Kazarovets, E. V., Durlevich, O. V., et al. (2017).
General catalogue of variable stars: Version GCVS 5.1.
\textit{Astronomy Reports}, 61(1), 80--88.

\bibitem[Sandler et al.(2018)]{sandler2018mobilenetv2}
Sandler, M., Howard, A., Zhu, M., et al. (2018).
MobileNetV2: Inverted Residuals and Linear Bottlenecks.
In \textit{Proceedings of the IEEE Conference on Computer Vision and Pattern Recognition} (pp. 4510--4520).

\bibitem[Sarro et al.(2009)]{sarro2009automated}
Sarro, L. M., Debosscher, J., López, M., \, \ &  Aerts, C. (2009).
Automated supervised classification of variable stars in the CoRoT programme.
\textit{Astronomy \& Astrophysics}, 494(2), 739--768.

\bibitem[Scargle(1982)]{scargle1982studies}
Scargle, J. D. (1982).
Studies in astronomical time series analysis. II-Statistical aspects of spectral analysis of unevenly spaced data.
\textit{The Astrophysical Journal}, 263, 835--853.

\bibitem[Stassun et al.(2014)]{stassun2014empirical}
Stassun, K. G., Feiden, G. A., \, \ &  Torres, G. (2014).
Empirical calibration of the fine-structure constant.
\textit{New Astronomy Reviews}, 60, 1--28.

\bibitem[Szegedy et al.(2016)]{szegedy2016rethinking}
Szegedy, C., Vanhoucke, V., Ioffe, S., et al. (2016).
Rethinking the Inception Architecture for Computer Vision.
In \textit{Proceedings of the IEEE Conference on Computer Vision and Pattern Recognition} (pp. 2818--2826).

\bibitem[Torres et al.(2010)]{torres2010binary}
Torres, G., Andersen, J., \, \ &  Giménez, A. (2010).
Accurate masses and radii of normal stars: Modern results and applications.
\textit{Astronomy \& Astrophysics Review}, 18(1-2), 67--126.

\bibitem[VanderPlas(2018)]{vanderplas2018understanding}
VanderPlas, J. T. (2018).
Understanding the Lomb-Scargle Periodogram.
\textit{The Astrophysical Journal Supplement Series}, 236(1), 16.

\bibitem[Wang \& Oates(2020)]{wang2020cnn}
Wang, Z., \, \ &  Oates, T. (2020).
CNN-based classification of time series with spatial augmentation.
\textit{Neurocomputing}, 391, 153--161.

\bibitem[Aleo et al.(2023)]{aleo2023ztf}
Aleo, P. D., Miller, A. A., Bellm, E. C., et al. (2023).
The ZTF Source Classification Project: II. Period Stellar Variables.
\textit{arXiv preprint arXiv:2302.08491}.

\end{thebibliography}

\end{document}
